\section{从一个可以观察的问题开始}

给操作系统增加功能时，先写一句运行后能够判断真假的话。例如“支持磁盘”
过于宽泛；“给定合法描述符，ROM 在有限轮询次数内把 Stage 1 从 LBA 1
读到物理地址 \code{0x8000}，校验通过后跳转”就能指导实现。它指出了输入、
目的地址、终止条件和成功动作。

再写出最值得担心的失败。这个例子至少包含描述符损坏、LBA 越界、设备 ERR
和 BSY 超时。成功标记与四种错误标记确定以后，测试才知道应构造什么镜像，
调试者也知道最后一行串口输出意味着什么。

\topic{先找出新机制依赖的上一层}
内核功能之间有天然顺序。没有串口时，磁盘失败很难定位；没有稳定的 ELF
装载，C++ 内核入口无从验证；没有异常入口和页帧所有权，用户进程越多，故障
越难收拾。开始编码前先列出新机制将依赖的状态，并回到测试确认这些状态真的
存在。

\begin{osgridlongtable}{L{0.22\textwidth}L{0.29\textwidth}L{0.37\textwidth}}
\textbf{准备增加的机制} & \textbf{必须先具备} & \textbf{最先观察什么}\\ \hline
Stage 1 磁盘加载 & 串口、低端 RAM、PIO 端口 & 描述符校验与目标内存字节\\ \hline
长模式 & GDT、A20、页表和有效栈 & 每次模式切换后的第一条日志\\ \hline
用户进程 & IDT/TSS、独立 CR3、用户复制 & Ring 3 入口和受控退出\\ \hline
fork/COW & VMA、页帧引用、页故障恢复 & 父子写入后的字节隔离\\ \hline
journal & 块缓存、flush、可重放盘面格式 & 每个断电点后的重新挂载结果\\ \hline
\end{osgridlongtable}

\topic{把一次改动压到一个能讲清楚的范围}
同一个提交里同时改页表、调度器、文件格式和 Shell，任何失败都会出现四条
候选路径。更容易审查的做法是先让一个层次成立：先建立 VFS 和 memfs，确认
路径解析；再让 rootfs 成为另一个后端；最后增加持久格式。每一步都保留可以
启动的系统。

\topic{实验记录写现象和推理}
一次有用的记录包括复现命令、预期、实际最后一行日志、当前假设、检查动作和
结论。不要只保存“成功了”的截图。若串口停在 \code{PAE_READY}，记录下一步
检查 EFER、CR0.PG、远跳目标和当前映射；若检查否定了某个猜测，也把否定依据
留下。

\topic{运行以后先回答三个问题}
\begin{enumerate}[label=\arabic*.]
  \item CPU 或当前 Thread 最后执行到了哪里？
  \item 此前写入的寄存器、内存或磁盘字节与预期是否相同？
  \item 下一次状态变化由普通控制流、异常、中断还是另一个 Thread 触发？
\end{enumerate}

这三个问题能把“机器卡住了”缩小为具体边界。此时再增加一条低频日志、检查
QEMU 寄存器或反汇编目标字节，得到的信息会比到处打印更有用。
